% !TEX root = main.tex
\clearpage
\subsection{Экспериментальный стенд}

Контейнеры \code{node-a} и \code{node-b} имели доступ к одному набору из четырёх виртуальных процессоров (CPU 0--3). Для MPI использовался транспорт TCP в изолированной виртуальной сети.

Сборка выполнена на \code{node-a} компилятором GNU C++~13.3.0 с ключами \code{-O3} и \code{-std=c++17}. Использованы GNU OpenMP и Open MPI~4.1.6~\cite{ompi}. Последовательный и OpenMP-варианты запущены на \code{node-a}. В двух остальных вариантах на каждом узле работал один MPI-процесс. На второй узел скопирован файл, собранный на первом. Паспорт стенда приведён в \code{results/environment.txt}.

\subsection{Методика измерений}

Для каждого размера ($257$, $512$, $1024$, $2048$) в каждом режиме выполнены один прогрев и семь измерений. В последовательном и OpenMP-вариантах через \code{steady_clock} измерялось время ядра. В режимах с MPI после общего барьера измерялось время рассылки, вычислений и сбора результата. Через \code{MPI_Reduce} на ранге 0 выбиралось наибольшее время среди рангов (листинг~\ref{lst:mpi}).

Измеряемый интервал начинается после выделения массивов, генерации данных и запуска процессов. Результат последовательного прогрева сохранён как эталон. После остановки таймера результат сравнивался с эталоном поэлементно. При несовпадении программа сообщает координаты первого различия.

Для каждой серии найдена медиана времени $T$. Ускорение $S$ и эффективность $E$ определены по формулам
\begin{equation}
S=\frac{T_{\text{посл}}}{T},\qquad E=\frac{S}{p},
\label{eq:metrics}
\end{equation}
где $T_{\text{посл}}$~--- медиана последовательного варианта при том же $N$, а $p$~--- суммарное число вычислительных потоков по всем процессам. В таблице~\ref{tab:configs} для OpenMP, MPI и гибридного режима этому соответствуют 4, 2 и 4 потока. Минимум, медиана и максимум каждой серии приведены в \code{results/summary.csv}.
